% Section 16.10, from lectures/L16/appendix/b_exercises.md.

\section{Exercises}
\label{c16:sec:exercises}

\begin{aside}
You are building the transmit path, and \code{can_controller_tb} drives two nodes through a complete
exchange, so it cannot pass until \chapref{c17:ch} adds the receive half. \code{make build-project}
knows that and reports it as \emph{skipped} instead of running it and failing:

\begin{console}
==> can_controller_tb (skipped: can_controller.vhd has no receive path yet; ...)
\end{console}

So the build stays green this chapter, and these exercises reason their way through the transmit path
and analyse it in GHDL instead. The testbench becomes the verdict one chapter from now.
\end{aside}

\exercise{Trace a short frame by hand}{Reasoning}
\label{c16:ex:shortframe}
From the frame format in \secref{c10:sec:frameformat} and this chapter's field table, write out the
whole sequence of states \code{can_controller} passes through (in the TX role) to send a frame with ID
\code{0x001}, DLC \code{0} (no data field), from \code{STATE_IDLE} back to \code{STATE_IDLE}. Note for
every state roughly how many bit periods it occupies.

\exercise{Why the \code{STATE_LOAD_*} state exists}{Reasoning}
\label{c16:ex:loadstate}
\xpart{a} Explain in your own words what would go wrong if a stuffed field were sequenced with a
single state that asserted \code{tx_shift_reg}'s \code{load} \emph{and} waited for \code{done},
instead of a separate one-cycle \code{STATE_LOAD_*} followed by the shift state.

\xpart{b} Point out the specific timing rule from \chapref{c14:ch} (that \code{tx_shift_reg}'s
\code{load} presents immediately, and its \code{done} comes one extra shift later) that makes the
split into two states necessary.

\exercise{What reaches \code{crc15}, and what does not}{Reasoning}
\label{c16:ex:crcfeed}
\xpart{a} For a transmitted frame, sort these categories of bits into ``fed into \code{crc15}'' and
``excluded'': the SOF bit, real bits from the ID, control field and data, inserted stuff bits, the CRC
field's own bits, and the unstuffed tail (CRC delimiter, ACK slot, ACK delimiter, EOF). For every
excluded category, say what excludes it: either a term in
\code{crc_enable <= txsr_bit_valid and not txsr_stuff and role;}, or the fact that the field never
passes through \code{tx_shift_reg} at all.

\xpart{b} One of those categories may surprise you: the CRC field's own bits go back into the engine.
Using \chapref{c13:ch}'s property that generating and checking are the same operation, say what value
the register holds after the last CRC bit has been fed in, and name the check in \chapref{c17:ch} that
depends on precisely that value. Then explain why the transmit path, taken on its own, does not depend
on it: what does \code{STATE_START}'s \code{clear} already guarantee about the second frame a node
sends?

\xpart{c} Suppose you had instead gated \code{crc_enable} off during \code{STATE_CRC_HI} and
\code{STATE_CRC_LO}. Would the first frame a node sends carry a correct CRC? Would the second? Be
precise: that \code{crc15} is cleared is part of the answer. Then say whether any simulation in this
book would catch the change, and what that tells you about where the fed-back CRC field actually
earns its keep.

\exercise{How long is a frame, exactly?}{Reasoning}
\label{c16:ex:framelen}
Answering that requires holding the whole field table and the bit stuffing rule in your head at once,
which is the point. Take the frame \code{can_controller} sends for ID \code{0x123}, DLC \code{3}, data
bytes \code{FF 00 00}.

\xpart{a} Write out the stuffed region bit by bit, SOF through the end of the data field, using this
chapter's field table. (It should come to 43 bits: $1 + 11 + 1 + 2 + 4 + 24$.) The identifier is
\code{0x123}; remember that RTR, IDE and r0 are all \code{0} and that DLC is \code{0011}.

\xpart{b} The CRC-15 over those 43 bits is \code{0x0FEE}, that is, \code{000111111101110} (you can
confirm that against your \code{crc15} from \chapref{c13:ch} if you like). Add it, which gives the
whole stuffed region of 58 bits, SOF through CRC.

\xpart{c} Apply the bit stuffing rule over all 58 bits and count the inserted stuff bits. Watch the
boundary between the data field's closing zeros and the CRC's leading zeros: the run counter is not
cleared at field boundaries.

\xpart{d} Add the unstuffed tail (CRC delimiter, ACK slot, ACK delimiter, EOF) and state the total
number of bits actually driven out on the wire. How many microseconds is that at 1~Mbit/s?

\xpart{e} The same frame with the data \code{11 22 33} needs only one stuff bit instead of the number
you got in \textbf{c)}. Explain in one sentence why a frame's duration on a CAN bus depends on the
\emph{payload's value} and not only on its length, and what that means for anyone trying to work out
worst-case bus timing.
